\section{Introduction}
\label{sec:introduction}

MQTT is a lightweight publish--subscribe protocol used across monitoring, industrial, transportation, home, and smart-city deployments. A recent Internet-wide measurement found 14,386 publicly reachable brokers and collected 3.2~billion messages, confirming both broad use and substantial diversity in topic and traffic behavior~\cite{innamorati2026pulse}. Edge and gateway applications increasingly need to combine numeric telemetry with operator-written alarm descriptions, heterogeneous vendor logs, maintenance histories, and natural-language explanations. An LLM is useful when these inputs cannot be reduced to one fixed parser or threshold table, while exact action labels remain valuable for reproducible safety evaluation. This combination creates an end-to-end control problem: protocol delivery may be fast, while the application invoked by a delivered request may queue for a heterogeneous CPU/GPU resource and run for orders of magnitude longer.

MQTT delivery and application completion have distinct semantics. In MQTT QoS~1, PUBACK transfers protocol responsibility for an Application Message. It does not report that the receiving subscriber has completed the action encoded by that message; indeed, MQTT does not require onward delivery to finish before PUBACK is sent~\cite{oasis2019mqtt}. MQTT~5 supplies a request/response pattern through Response Topic and Correlation Data, but the broker treats the request and response as ordinary Application Messages. Message Expiry Interval prevents stale undelivered messages from being forwarded indefinitely; it does not retain a delivered request until LLM execution finishes. Consequently, completion-aware retention requires application state separate from standard acknowledgment state.

We address this problem with the \emph{Deadline- and Quality-Aware Application-state Retention Controller} (\method{}), a distinct bounded controller co-located with the broker. A publisher sends a normal MQTT~5 request. The controller admits a record if capacity remains, forwards the job with an application deadline, and releases the record on the first correlated completion or expiry. This record store is the controlled resource. Rejection at this store is termed \emph{application-state blocking} and is distinct from MQTT packet loss and native in-flight blocking.

The second difficulty is model mismatch. Erlang--B gives an interpretable loss budget under locally stationary Poisson arrivals and independent holding times, but an exact inversion of that formula cannot itself demonstrate that a deployed system meets a target. IoT aggregates can be approximately Poisson only under scale, independence, and synchronization conditions~\cite{metzger2019modeling,lopez2020aggregated}; alarm storms and event-triggered sources can be bursty. LLM completion time is also configuration-, hardware-, queue-, and output-dependent. \method{} uses Erlang--B as a feedforward component, converts its mean holding-time budget into a deadline using measured completion samples, and corrects the deadline with observed blocking. A guard adds class-specific minimum budgets and load-aware routing when the common loss-model deadline would make useful completion infeasible.

Finally, longer prompts and outputs are not synonymous with higher task utility. Long-context studies report performance degradation with distractors and position effects~\cite{liu2024lost,hsieh2024ruler}, while LLM-serving systems explicitly separate prefill and decode costs and optimize latency/goodput rather than token count alone~\cite{kwon2023vllm,agrawal2024sarathi,zhong2024distserve}. Our subscriber therefore selects a measured configuration by held-out strict task accuracy subject to its tail-latency budget. The context limit is an inequality; unused context and early completion are valid outcomes.

The contributions are as follows.

\begin{itemize}
  \item We define and implement a protocol-correct MQTT~5 request/completion path that separates PUBACK, Message Expiry, and correlated application completion, and controls a bounded application record store.
  \item We develop a measured controller that combines conditional Erlang--B feedforward, empirical truncated-mean inversion, blocking feedback, class-specific feasibility floors, and load-aware rerouting across heterogeneous execution classes.
  \item We formulate inference configuration as a finite class-conditioned selection problem using calibration correctness and profiled latency. The 240-item corpus includes direct readings, temporal histories, and heterogeneous logs, so different requests can require different context.
  \item We evaluate the complete loop using a real Mosquitto broker, three Qwen2.5 sizes, and a Llama-3.1--8B-class comparator. Ten paired replications per condition compare \method{} with occupancy-threshold admission, static class reservation, and their combination under Poisson, MMPP, and simultaneous arrival-rate/model-service drift, while matched studies cover four allocation rules and a five-point fixed-deadline sweep.
  \item We characterize the operating region under model mismatch and burstiness. Empirical replay includes a bounded-delay token-bucket variant and reports shaper loss separately from broker blocking, while an explicit local stability condition states when the blocking-feedback subsystem converges.
\end{itemize}

The remainder of the paper reviews related work, defines the protocol and controlled state, presents the controller and subscriber policy, evaluates the implementation, and discusses limitations and deployment implications.
